\documentclass[aspectratio=169,11pt]{beamer}
\usetheme{Madrid}
\usecolortheme{default}
\setbeamertemplate{navigation symbols}{}
\setbeamertemplate{footline}[frame number]
\usepackage{amsmath,amssymb,mathtools}
\usepackage{graphicx}
\newcommand{\Smat}{\mathbf S}
\newcommand{\fvec}{\mathbf f}
\newcommand{\yvec}{\mathbf y}
\newcommand{\evec}{\mathbf e}
\newcommand{\Rmat}{\mathbf R}
\newcommand{\pvec}{\mathbf p}
\title{Least-Squares Filter Design}
\subtitle{From prediction error to the normal equations}
\author{ERTH 467 --- Geoscientific Signal Processing}
\date{}
\begin{document}
\begin{frame}\titlepage\end{frame}

\begin{frame}{What are we trying to achieve?}
\begin{itemize}
\item Known input signal $s_n$ and desired output (target) $y_n$.
\item Design a filter $f_k$ so that the convolution predicts the target:
\end{itemize}
\[\widehat y_n=\sum_{k=0}^{M-1}f_k s_{n-k}\]
\begin{block}{Matrix representation}
\[\boxed{\widehat{\yvec}=\Smat\fvec}\]
\end{block}
\small $\Smat$ contains shifted versions of the input signal; $\fvec$ contains the unknown coefficients.
\end{frame}

\begin{frame}{Convolution as a matrix multiplication}
\small For zero-padded causal convolution, a Toeplitz matrix is typical:
\[\underbrace{\begin{bmatrix}s_0&0&0\\s_1&s_0&0\\s_2&s_1&s_0\\s_3&s_2&s_1\\\vdots&\vdots&\vdots\end{bmatrix}}_{\Smat}\underbrace{\begin{bmatrix}f_0\\f_1\\f_2\end{bmatrix}}_{\fvec}=\underbrace{\begin{bmatrix}\widehat y_0\\\widehat y_1\\\widehat y_2\\\widehat y_3\\\vdots\end{bmatrix}}_{\widehat\yvec}\]
\begin{itemize}
\item $\Smat$: $N\times M$, $\fvec$: $M\times1$, target $\yvec$: $N\times1$.
\item \textbf{Periodic convolution} gives a \emph{circulant} $\Smat$ instead.
\item The following least-squares derivation works for \emph{either} matrix.
\end{itemize}
\end{frame}

\begin{frame}{Define the prediction error}
\[\boxed{\evec=\yvec-\widehat\yvec=\yvec-\Smat\fvec}\]
\begin{itemize}
\item Usually we cannot match every target sample exactly.
\item We therefore choose the filter that minimizes the \textbf{sum of squared errors}.
\end{itemize}
\begin{block}{Least-squares objective: a scalar function of $\fvec$}
\[\boxed{J(\fvec)=\sum_{n=0}^{N-1}e_n^2=\evec^T\evec=(\yvec-\Smat\fvec)^T(\yvec-\Smat\fvec)}\]
\end{block}
\end{frame}

\begin{frame}{How do we locate the minimum?}
\begin{itemize}
\item In elementary calculus: set the derivative of a differentiable function to zero.
\item Here $J$ depends on $M$ filter coefficients $f_0,\ldots,f_{M-1}$.
\item Set every partial derivative to zero, or equivalently the gradient:
\end{itemize}
\[\nabla_{\fvec}J=\begin{bmatrix}\partial J/\partial f_0\\\partial J/\partial f_1\\\vdots\\\partial J/\partial f_{M-1}\end{bmatrix}=\mathbf0\]
\begin{alertblock}{Our task}
Compute $\nabla_{\fvec}J$ from the squared-error expression.
\end{alertblock}
\end{frame}

\begin{frame}{Step 1: expand the squared prediction error}
\small
\begin{align*}
J(\fvec)&=(\yvec-\Smat\fvec)^T(\yvec-\Smat\fvec)\\[4pt]
&=(\yvec^T-\fvec^T\Smat^T)(\yvec-\Smat\fvec)\\[4pt]
&=\yvec^T\yvec-\yvec^T\Smat\fvec-\fvec^T\Smat^T\yvec+\fvec^T\Smat^T\Smat\fvec.
\end{align*}
\begin{block}{Key observation}
$\yvec^T\Smat\fvec$ and $\fvec^T\Smat^T\yvec$ are equal scalars (transposes of each other).
\end{block}
\[\boxed{J(\fvec)=\yvec^T\yvec-2\fvec^T\Smat^T\yvec+\fvec^T\Smat^T\Smat\fvec}\]
\end{frame}

\begin{frame}{Step 2: differentiate term by term}
\small Only $\fvec$ varies; $\Smat$ and $\yvec$ are known.
\begin{align*}
\nabla_{\fvec}(\yvec^T\yvec)&=\mathbf0&&\text{constant term}\\[7pt]
\nabla_{\fvec}(-2\fvec^T\Smat^T\yvec)&=-2\Smat^T\yvec&&\text{linear term}\\[7pt]
\nabla_{\fvec}(\fvec^T\Smat^T\Smat\fvec)&=2\Smat^T\Smat\fvec&&\text{quadratic term}
\end{align*}
\vspace{3mm}
\begin{block}{Why a factor of 2 in the last term?}
$\Smat^T\Smat$ is symmetric; for symmetric $\mathbf A$, $\nabla_{\fvec}(\fvec^T\mathbf A\fvec)=2\mathbf A\fvec$.
\end{block}
\end{frame}

\begin{frame}{Step 2: collect the gradient}
\[\boxed{\nabla_{\fvec}J=-2\Smat^T\yvec+2\Smat^T\Smat\fvec}\]
\vspace{6mm}
Factor the expression:
\[\boxed{\nabla_{\fvec}J=-2\Smat^T(\yvec-\Smat\fvec)=-2\Smat^T\evec}\]
\vspace{5mm}
\begin{itemize}
\item Each gradient component describes the change in total squared error when one filter coefficient changes.
\item We seek a filter for which \textbf{all} these derivatives vanish.
\end{itemize}
\end{frame}

\begin{frame}{Step 3: set the gradient to zero}
\begin{align*}
\nabla_{\fvec}J&=\mathbf0\\[5pt]
-2\Smat^T\yvec+2\Smat^T\Smat\fvec&=\mathbf0\\[5pt]
\Smat^T\Smat\fvec&=\Smat^T\yvec
\end{align*}
\begin{block}{Normal equations}
\[\boxed{\Smat^T\Smat\fvec=\Smat^T\yvec}\]
A system of $M$ linear equations for the $M$ unknown filter coefficients.
\end{block}
\end{frame}

\begin{frame}{Solving for the least-squares filter}
\begin{block}{Formal solution (when $\Smat^T\Smat$ is invertible)}
\[\boxed{\fvec_{\mathrm{LS}}=(\Smat^T\Smat)^{-1}\Smat^T\yvec}\]
\end{block}
\begin{itemize}
\item This is the expression shown on the lecture slide.
\item \textbf{Numerically:} solve $\Smat^T\Smat\fvec=\Smat^T\yvec$; do not explicitly form the inverse.
\item If $\Smat$ has dependent columns, an ordinary inverse does not exist; a pseudoinverse or regularization can be used.
\end{itemize}
\end{frame}

\begin{frame}{Why is this stationary point a minimum?}
\small The Hessian (matrix of second derivatives) is
\[\nabla_{\fvec}^2J=2\Smat^T\Smat.\]
For any vector $\mathbf v$,
\[\mathbf v^T\Smat^T\Smat\mathbf v=\|\Smat\mathbf v\|_2^2\ge0.\]
\begin{block}{Consequence}
$J$ is a \textbf{convex quadratic}: any solution of the normal equations is a global minimum. If $\Smat$ has independent columns, this minimum is unique.
\end{block}
\end{frame}

\begin{frame}{Geometric meaning: the residual is orthogonal}
At the optimal filter,
\[\nabla_{\fvec}J=-2\Smat^T\evec_{\mathrm{LS}}=\mathbf0\]
so
\[\boxed{\Smat^T\evec_{\mathrm{LS}}=\mathbf0}.\]
\begin{itemize}
\item The remaining prediction error is orthogonal to every column of $\Smat$.
\item No further adjustment using the available shifted input signals can reduce the squared error.
\end{itemize}
\end{frame}

\begin{frame}{Connection to correlation}
Define
\[\Rmat=\Smat^T\Smat,\qquad \pvec=\Smat^T\yvec.\]
Then the normal equations read
\[\boxed{\Rmat\fvec=\pvec}.\]
\begin{itemize}
\item $\Rmat$: inner products between shifted copies of the \textbf{input} (autocorrelation-related).
\item $\pvec$: inner products between shifted input and \textbf{target} (cross-correlation-related).
\item Boundary conventions (zero padding vs periodicity) affect exact entries.
\end{itemize}
\end{frame}

\begin{frame}{Optional: derive the equations one coefficient at a time}
\scriptsize
\[J=\sum_{n=0}^{N-1}\left(y_n-\sum_{k=0}^{M-1}S_{nk}f_k\right)^2\]
By the chain rule, differentiating with respect to a single $f_j$ gives
\begin{align*}
\frac{\partial J}{\partial f_j}
&=2\sum_{n=0}^{N-1}\left(y_n-\sum_{k=0}^{M-1}S_{nk}f_k\right)(-S_{nj})\\
&=-2\sum_{n=0}^{N-1}S_{nj}e_n.
\end{align*}
Setting this equal to zero and rearranging:
\[\boxed{\sum_{n=0}^{N-1}S_{nj}\sum_{k=0}^{M-1}S_{nk}f_k=\sum_{n=0}^{N-1}S_{nj}y_n}\]
This is the $j$th row of $\Smat^T\Smat\fvec=\Smat^T\yvec$.
\end{frame}

\begin{frame}{Summary: from prediction error to filter}
\begin{enumerate}
\item Predict: $\widehat\yvec=\Smat\fvec$.
\item Define error: $\evec=\yvec-\Smat\fvec$.
\item Minimize $J(\fvec)=\evec^T\evec$.
\item Differentiate: $\nabla_{\fvec}J=-2\Smat^T\yvec+2\Smat^T\Smat\fvec$.
\item Set equal to zero and solve:
\end{enumerate}
\vspace{5mm}
\begin{block}{Main result}
\[\boxed{\Smat^T\Smat\fvec=\Smat^T\yvec}\]
\end{block}
\end{frame}
\end{document}
